Pusat Data dan Informasi (Pusdatin) Universitas Tarumanagara (UNTAR) merupakan unit pengelola infrastruktur teknologi informasi dan sistem data terpusat yang mendukung seluruh kegiatan operasional akademik maupun administratif di lingkungan universitas. Seiring berjalannya waktu dan terus meningkatnya jumlah mahasiswa yang terdaftar serta transaksi akademik yang terjadi setiap semester, volume data yang tersimpan di dalam basis data produksi (\textit{active database}) mengalami pertumbuhan yang sangat signifikan dan pesat. Akumulasi data historis ini mencakup rekam jejak akademik, transkrip nilai, data kehadiran, hasil evaluasi pembelajaran, hingga log transaksi operasional mahasiswa dari seluruh angkatan tahun-tahun sebelumnya.

Penumpukan data historis mahasiswa dengan volume yang sangat besar dalam satu basis data produksi utama menimbulkan beberapa kendala krusial yang berdampak langsung pada operasional harian Pusdatin UNTAR. Pertama, kapasitas ruang penyimpanan (\textit{storage space}) pada basis data produksi kian terbatas dan terbebankan secara tidak proporsional oleh data historis yang sudah jarang diakses secara langsung dalam kegiatan operasional harian. Kedua, akumulasi tabel yang terus membengkak berdampak langsung pada penurunan efisiensi dan performa eksekusi kueri (\textit{query performance}) untuk layanan transaksi aktif mahasiswa, dosen, maupun staf administrasi. Di sisi lain, data historis mahasiswa tersebut sama sekali tidak dapat dihapus secara permanen (\textit{hard delete}) karena terikat oleh kebutuhan regulatori, kebijakan retensi dokumen akademik universitas, serta kewajiban pemenuhan kepatuhan audit internal maupun eksternal yang bersifat jangka panjang \cite{zulkifli2011manajemen}.

Selain masalah performa kueri dan keterbatasan ruang penyimpanan, kendala signifikan juga ditemukan pada proses pencadangan data (\textit{data backup}) yang berjalan saat ini di Pusdatin UNTAR. Proses pencadangan dilakukan secara menyeluruh (\textit{full database backup}) yang mencakup gabungan antara data transaksi aktif dan data historis secara terkonsolidasi pada satu server utama. Pendekatan ini membutuhkan durasi waktu yang sangat lama (\textit{time-consuming}) serta memerlukan alokasi sumber daya komputasi yang tinggi, sehingga dinilai kurang efisien dan berpotensi meningkatkan risiko kegagalan proses pencadangan pada saat terjadi lonjakan transaksi operasional.

Lebih lanjut, ketika unit terkait seperti Biro Administrasi Akademik atau tim auditor memerlukan data historis mahasiswa tertentu untuk keperluan pelaporan maupun pemeriksaan, tidak terdapat mekanisme yang terstruktur dan cepat untuk menemukannya. Pencarian manual yang dilakukan langsung pada basis data produksi berisiko mengganggu kinerja sistem transaksi aktif yang sedang berjalan, sekaligus berpotensi mengekspos integritas data operasional terhadap perubahan yang tidak terotorisasi \cite{james2014database}.

Guna mengatasi permasalahan tersebut secara sistematis, diperlukan perancangan mekanisme pengarsipan data (\textit{data archiving}) otomatis yang mampu memisahkan data historis mahasiswa dari basis data produksi ke dalam basis data arsip yang terpisah secara terjadwal dan terkonfigurasi. Sistem pengarsipan ini dirancang dengan arsitektur \textit{Client-Worker} terdistribusi yang berkomunikasi melalui HTTP API JSON (\textit{agent.ashx}), di mana eksekusi antrean tugas (\textit{task queue}) dijalankan oleh SQL Server Agent dan skrip PowerShell pada mesin eksternal terpisah. Dengan arsitektur ini, beban komputasi proses pengarsipan tidak langsung menimpa infrastruktur server kampus yang kapasitasnya terbatas. Sistem ini juga dilengkapi dengan portal \textit{on-demand retrieval} berbasis \textit{web} ASP.NET yang memungkinkan staf dan unit terkait di Pusdatin UNTAR melakukan pencarian dan penarikan data historis secara cepat dan akurat lintas basis data---baik pada basis data aktif maupun basis data arsip---tanpa membebani basis data produksi utama.

Berdasarkan uraian permasalahan dan gagasan solusi di atas, penulis mengangkat penelitian skripsi dengan judul \textbf{"PERANCANGAN APLIKASI DATA \textit{ARCHIVING} OTOMATIS DAN PORTAL \textit{ON-DEMAND RETRIEVAL} BERBASIS WEB ASP.NET PADA PUSDATIN UNTAR"}. Penelitian ini diharapkan dapat membantu Pusdatin UNTAR dalam meningkatkan efisiensi penyimpanan data, mengoptimalkan kecepatan proses pencadangan, serta mempercepat akses penarikan data historis mahasiswa sesuai kebutuhan audit, regulasi, dan pelayanan akademik yang berlaku.
